\documentclass[]{article}
\usepackage[spanish.mexico]{babel}
\usepackage[T1]{fontenc}
\usepackage[utf8]{inputenc}
%\usepackageB.\\{lmodern}
\usepackage[a4paper]{geometry}
 \geometry{
	a4paper,
	total={170mm,257mm},
	left=20mm,
	top=20mm,
}

%columnas itemize
\usepackage{multicol}

%\usepackage{natbib}
\usepackage{cite}


%Grafico de barras
\usepackage{pgfplots}

%Graficos e imagenes
\usepackage{graphicx}


\usepackage{tikz}
\usepackage[american voltages, american currents,siunitx]{circuitikz}


\title{Examen Máquinas Eléctricas}
\author{Pablo Vivar Colina}
%\date{Mayo 2018}



\begin{document}
	
%	%\usepackage[top=2cm,bottom=2cm,left=1cm,right=1cm]{geometry}


\begin{titlepage}
     \begin{center}
	\includegraphics[width=0.09\textwidth]{UNAM}\Large Universidad Nacional Autónoma de México
        	\includegraphics[width=0.09\textwidth]{FI}\\[1cm]
        \Large Facultad de Ingeniería\\[1cm]
       % \Large División de Ciencias Básicas\\[1cm]
         \Large Laboratorio de Fundamentos de Control(6655)\\[1cm]
         %la clave antes era:4314
         \footnotesize Profesor: Salcedo Ubilla María Leonor Ing.\\[1cm]
        \footnotesize Semestre 2019-1\\[1cm]
        
       

        \Large Práctica No. 1\\[1cm]
        
           

\Large Introdcción MATLAB
        
         %Texto a la derecha
          \begin{flushright}
\footnotesize  Grupo 2\\[0.5cm]
\footnotesize Brigada: 4\\[0.5cm]
\footnotesize Rodrigo Adrián Martínez López\\[0.5cm]
\footnotesize Vivar Colina Pablo\\[0.5cm]
 \end{flushright}
    %Texto a la izquierda
          \begin{flushleft}
        \footnotesize Ciudad Universitaria Agosto de 2018.\\
          \end{flushleft}
         
          
        %\vfill
        %\today
   \end{center}
\end{titlepage}
 %agregar portada

\maketitle

%\tableofcontents  % Write out the Table of Contents

%\listoffigures  % Write out the List of Figures



\section{Transformador}

\begin{multicols}{2}
	\begin{itemize}
		\item P=[15 KVA]
		\item V=2400[V] a 240 [V]
		\item f=60 [Hz]
		\item $R_2$=0.025 [$\Omega$]
		\item $R_C$=32[k$\Omega$]
		\item $X_m$=11.5 [k$\Omega$]
		\item $X_1$=7[$\Omega$]
		\item $X_2$=0.07[$\Omega$]
	\end{itemize}
\end{multicols}

\section{Motor trifásico}

\begin{multicols}{2}
	\begin{itemize}
		\item V=230 [V]
		\item P=5[Hp]=3.7282[W]
		\item f=60 [Hz]
		\item V=1746 [rpm]
		\item I=12[A]
		\item 2.5 $\%$ rotor
		\item 2.5 $\%$ estator
		\item 5 $\%$ núcleo
		\item 3 $\%$ devanado
		\item $X_1$=2.7 [$\Omega$]
		\item $X_R$=37 [$Omega$]
	\end{itemize}

\end{multicols}

\begin{figure}[h!]
	\centering
%	\begin{center}
		\begin{circuitikz}
			
			\draw
			
			
			(0,0)to[R,l=$R_1$25](2,0)
			(2,0)to[L,l=$L_1$7](4,0)
			
			
			(4,0)to[R,l=$R_2$0.025](6,0)
			(6,0)to[L,l=$L_2$11.5](8,0)
			
			(4,0)--(4,-1)
			(3,-1)--(5,-1)
			
			(3,-1)to[R,l=$R_3$0.025](3,-3)
			(5,-1)to[L,l=$L_3$0.007](5,-3)
			
			(3,-3)--(5,-3)
			
			(4,-3)--(4,-4)
			
			(0,-4)--(8,-4)
			
			(8,-4)--(8,0)
			
			%(8,-2)node[anchor=east] {240V}
			(8,-2)node[anchor=west] {240V}
			
			
			(8,-3)rectangle(9,-1)
			
			(9,-4)--(9,0)
			
			%LO MISMO MAS 9
			
			(0+9,0)to[R,l=$R_4$1K](2+9,0)
			(2+9,0)to[L,l=$L_4$27](4+9,0)
			
			
			(4+9,0)to[R,l=$R_5$1K](6+9,0)
			(6+9,0)to[L,l=$L_5$1m](8+9,0)
			
			(4+9,0)--(4+9,-1)
			(3+9,-1)--(5+9,-1)
			
			(3+9,-1)to[R,l=$R_6$1K](3+9,-3)
			(5+9,-1)to[L,l=$L_6$0](5+9,-3)
			
			(3+9,-3)--(5+9,-3)
			
			(4+9,-3)--(4+9,-4)
			
			(9,-4)--(8+9,-4)
			
			(8+9,-4)--(8+9,0)


			
			%(8,-2)node[anchor=east] {240V}
			(8+9,-2)node[anchor=west] {$\frac{1-s}{s}$}
				
          (8+9,-3)rectangle(9+9,-1)
			
		
			
			%tierras
			(4,-3) to   (4,-4) node[ground]{}
			(4+9,-3) to   (4+9,-4) node[ground]{}
			
			
			
			
			%node[scale=1.2,nand port,anchor=in 1] (nand1) {}
			%(nand1.in 2) -- (xor1.in 2|-nand1.in 2)
			
			%(3,0)to[push button,l=$R_1$](3,3)
			
			
			%Generador de funciones
			%(-6,0.5)--(-5,0.5)
			%(0,3)to[V,l=$G$](0,0)
			
			;
			
		\end{circuitikz}
		
%\end{center}
	
		
%\caption{Circuito de acción momentánea}
%\label{fig:CircuitoMomentáneo}
\end{figure}

Determinar

\begin{enumerate}
	\item Velocidad Síncrona
	\item Deslizamiento
	\item La frecuencia de las corrientes del motor
	\item Par de Salida
	\item Corriente de suministro del motor
	\item La eficiencia del motor
	\item Factor de potencia
\end{enumerate}

\section{Desarrollo}

\subsection{Velocidad Síncrona}


\begin{equation}
  n_s=\frac{120(60)}{8}(1-0.03)=900
  \label{ns}
\end{equation}

\subsection{Deslizamiento}

%Del resultado de la ecuación \ref{ns} sustituimos:\\

\begin{equation}
s=\frac{n_s-n_m}{n_s}=\frac{1800-1746}{1800}=0.03
\end{equation}

\subsection{Frecuencia de las corrientes del motor}

\begin{equation}
f_i=sf=(0.03)(60)=1.8[Hz]
\end{equation}

\subsection{Par de Salida}

\begin{equation}
 \tau=\frac{P}{v}=\frac{3.7282[kW]}{1746[rpm](\frac{1[M]}{60[s]})2 \pi}=\frac{3.7282[kW]}{182[\frac{rad}{s}]}=20[Nm]
\end{equation}

\subsection{Corriente de suministro del motor}

\begin{equation}
\eta=\frac{P_{salida}}{P_{entrada}}=\frac{3728}{4285}=0.87=87\%
\end{equation}

\subsection{Factor de potencia}


\begin{equation}
  F.P.=\frac{4285}{(230)(8)\sqrt{3}}=0.59  
\end{equation}

\end{document}
